\documentclass[11pt,a4paper]{article}
\usepackage[margin=1in]{geometry}
\usepackage{amsmath,amssymb,amsthm,mathtools}
\usepackage[T1]{fontenc}
\usepackage{lmodern}
\usepackage{xcolor}
\usepackage[hidelinks]{hyperref}
\hypersetup{
  pdftitle={Fat Gauss Fibres and Tjurina--Milnor Defects Forced by Osculating Absorption},
  pdfauthor={Lluis Eriksson}
}

\newtheorem{theorem}{Theorem}[section]
\newtheorem{proposition}[theorem]{Proposition}
\newtheorem{lemma}[theorem]{Lemma}
\newtheorem{corollary}[theorem]{Corollary}
\theoremstyle{definition}
\newtheorem{definition}[theorem]{Definition}
\newtheorem{remark}[theorem]{Remark}

\newcommand{\kk}{\mathbb C}
\newcommand{\PP}{\mathbb P}
\newcommand{\cO}{\mathcal O}
\newcommand{\cI}{\mathcal I}
\newcommand{\m}{\mathfrak m}
\newcommand{\Osc}{\operatorname{Osc}}
\newcommand{\mult}{\operatorname{mult}}
\newcommand{\Sing}{\operatorname{Sing}}
\newcommand{\Fitt}{\operatorname{Fitt}}
\newcommand{\length}{\operatorname{length}}

\title{Fat Gauss Fibres and Tjurina--Milnor Defects\\
Forced by Osculating Absorption}
\author{Lluis Eriksson\\\small Independent researcher}
\date{August 25, 2026\quad (version 1.0)}

\begin{document}
\maketitle

\begin{abstract}
Let $X^d\subset\PP^{d+1}_{\mathbb C}$ be a smooth non-linear hypersurface,
let $\gamma_X:X\to X^\vee$ be its Gauss normalization, and let $\Gamma_\eta$
be the scheme-theoretic fibre over a tangent hyperplane $\eta=[W]$.
If $h_p$ is the isolated germ of the tangent section $X\cap W$ at a point
$p$ of the fibre, the first-jet description of the conormal incidence gives
\[
 \widehat{\cO}_{\Gamma_\eta,p}
 \cong \kk[[z_1,\ldots,z_d]]/(h_p,\partial h_p).
\]
Thus the fibre length is $\sum_p\tau(h_p)$, whereas the classical formula
of Dimca gives $\mult_\eta(X^\vee)=\sum_p\mu(h_p)$.  We record the exact
defect identity
\[
 \mult_\eta(X^\vee)-\length(\Gamma_\eta)
 =\sum_p(\mu_p-\tau_p).
\]
We also compute
$\Fitt_0\Omega_{X/X^\vee,p}=(\det\operatorname{Hess}h_p)$.
These local identities are classical or immediate consequences of classical
first-jet and singularity theory; they are not claimed as new.  The equality
$\mult_\eta(X^\vee)=\length(\Gamma_\eta)$ occurs precisely when all section
germs are analytically quasihomogeneous.

The application proved here is to point-span osculating absorption.  If the complete
$\cO_X(m)$ point span of the reduced fibre absorbs order $s$, $1\le s\le m$,
then $h_p\in\m_p^{s+1}$ and the Gauss fibre contains
$\operatorname{Spec}(\cO_{X,p}/\m_p^s)$ at every support.  Consequently
\[
 \length(\Gamma_\eta)
 \geq \binom{d+s-1}{d}|Z|
 \geq \binom{d+s-1}{d}\binom{d+m}{d},
\]
and the ramification equation vanishes to order at least $d(s-1)$ at each
support.  Finally, when $d(s-1)>s+1$, we construct smooth hypersurfaces for
which the reduced fibre has the extremal size $N=\binom{d+m}{d}$ and, for
every $0\le q\le N$,
\[
 \mult_\eta(X^\vee)=s^dN,
 \qquad \length(\Gamma_\eta)=s^dN-q.
\]
The construction uses finite contact determinacy and prescribed jets.  Its
degree threshold is sufficient rather than minimal, and no exhaustive
priority claim is made.
\end{abstract}

\section{Statement and scope}

Let $X^d\subset\PP^{d+1}_{\kk}$ be a smooth hypersurface of degree $D\ge2$.
Write
\[
 \gamma_X:X\longrightarrow X^\vee\subset(\PP^{d+1})^\vee
\]
for the Gauss morphism.  It is finite and birational, hence is the
normalization of the dual hypersurface; see, for example,
Tevelev~\cite[Chapter~1]{Tevelev2005} and the proof recalled in
\cite[Lemma~2.1]{ErikssonDual2026}.  For $\eta=[W]\in X^\vee$, put
\[
 \Gamma_\eta=X\times_{X^\vee}\operatorname{Spec}\kappa(\eta),
 \qquad Z=(\Gamma_\eta)_{\mathrm{red}}.
\]
All fibre schemes in this paper are fibres over the reduced complex point
$\eta$; no claim about a scheme structure on a varying contact locus is
implicit.

Fix $m\ge1$, put $H=\cO_X(1)$, and let
\[
 S_Z=\operatorname{Im}\!\left(
 H^0(Z,H^m|_Z)^*\longrightarrow H^0(X,H^m)^*\right).
\]
Principal-parts evaluation defines the affine order-$s$ osculating block
\[
 \widehat\Osc^s_p(H^m)=\operatorname{Im}\!\left(
 H^0((s+1)p,H^m|_{(s+1)p})^*
 \longrightarrow H^0(X,H^m)^*\right).
\]

\begin{definition}
For $1\le s\le m$, the reduced fibre $Z$ is
\emph{point-span $s$-osculating-absorbing} if
\[
 \widehat\Osc^s_p(H^m)\subseteq S_Z\qquad(p\in Z).
\]
\end{definition}

The first theorem separates three invariants that should not be conflated:
the cardinality of the reduced fibre, the length of the fibre scheme, and
the multiplicity of the dual.
For an isolated analytic hypersurface germ $h$, our conventions are
\[
 \mu(h)=\dim_\kk\frac{\kk\{z\}}{(\partial h)},
 \qquad
 \tau(h)=\dim_\kk\frac{\kk\{z\}}{(h,\partial h)}.
\]

\begin{theorem}[Gauss fibre, Tjurina length, and dual defect]\label{thm:local}
Let $\eta=[W]\in X^\vee$ and let $h_p\in\kk\{z_1,\ldots,z_d\}$ be a local
equation of the tangent section $X\cap W$ at $p\in Z$.  Then
\begin{align*}
 \widehat{\cO}_{\Gamma_\eta,p}
 &\cong \kk[[z_1,\ldots,z_d]]/(h_p,\partial_1h_p,\ldots,\partial_dh_p),\\
 \length(\Gamma_\eta)&=\sum_{p\in Z}\tau(h_p),\\
 \mult_\eta(X^\vee)&=\sum_{p\in Z}\mu(h_p).
\end{align*}
Consequently
\[
 \mult_\eta(X^\vee)-\length(\Gamma_\eta)
 =\sum_{p\in Z}\bigl(\mu(h_p)-\tau(h_p)\bigr)\ge0.       \tag{1}
\]
Equality in \emph{(1)} holds if and only if every $h_p$ is analytically
quasihomogeneous.  Moreover,
\[
 \Fitt_0\Omega_{X/X^\vee,p}
   =(\det\operatorname{Hess}h_p)\subset\cO_{X,p}^{\mathrm{an}}.       \tag{2}
\]
\end{theorem}

\begin{theorem}[Fat-fibre floor forced by absorption]\label{thm:floor}
Assume that $Z$ is point-span $s$-osculating-absorbing in the complete
$H^m$ embedding.  For every $p\in Z$ there is a closed immersion
\[
 \operatorname{Spec}(\cO_{X,p}/\m_{X,p}^s)\hookrightarrow\Gamma_\eta.
\]
In particular, with $N_{d,m}=\binom{d+m}{d}$,
\[
 \length(\Gamma_\eta)
 \ge \binom{d+s-1}{d}|Z|
 \ge \binom{d+s-1}{d}N_{d,m}.                         \tag{3}
\]
The local ramification equation in \emph{(2)} has order at least
$d(s-1)$ at every $p\in Z$.  If $s\ge2$, the fibre scheme is nonreduced at
every support.
\end{theorem}

The lower bound in (3) is not asserted to be optimal for $s>1$.  Our global
examples instead show that all integral defects from zero through the number
of extremal supports occur while every local Milnor number remains fixed at
$s^d$.

\begin{theorem}[Global realization of Tjurina--Milnor defects]\label{thm:realization}
Let $d,m,s$ be positive integers with $1\le s\le m$ and
\[
 d(s-1)>s+1.
\]
Put $N=N_{d,m}$ and suppose
\[
 D\ge (s^d+2)N+1.                                      \tag{4}
\]
For every integer $q$ with $0\le q\le N$, there exist a smooth integral
degree-$D$ hypersurface $X^d\subset\PP^{d+1}$, a hyperplane $W$, and
$\eta=[W]\in X^\vee$ such that:
\begin{enumerate}
\item $Z=(\gamma_X^{-1}(\eta))_{\mathrm{red}}$ has exactly $N$ points;
\item $\dim S_Z=N< h^0(X,H^m)$ and $Z$ is point-span
      $s$-osculating-absorbing;
\item every tangent-section germ has Milnor number $s^d$; exactly $q$ of
      them have Tjurina number $s^d-1$, and the others have Tjurina number
      $s^d$;
\item
\[
 \mult_\eta(X^\vee)=s^dN,
 \qquad \length(\Gamma_\eta)=s^dN-q;
\]
\item the ramification equation has order exactly $d(s-1)$ at every point
      of $Z$.
\end{enumerate}
The construction is existential up to nonempty Zariski-open choices.  The
threshold \emph{(4)} is not claimed to be minimal.
\end{theorem}

\section{The local fibre ideal}\label{sec:local}

Choose analytic coordinates centred at $p\in Z$ in which
\[
 X:\quad w=h(z_1,\ldots,z_d),
 \qquad W:\quad w=0,
 \qquad h\in\m^2.
\]
The tangent hyperplane at $(z,h(z))$ has equation
\[
 Y-\sum_i h_i(z)Z_i+\left(\sum_i z_i h_i(z)-h(z)\right)=0.
\]
On the dual chart where the coefficient of $Y$ is one, the Gauss map is
therefore, up to harmless signs,
\[
 z\longmapsto\left(h_1(z),\ldots,h_d(z),
 h(z)-\sum_i z_i h_i(z)\right).                         \tag{5}
\]
The extension of the maximal ideal of $\eta$ is
\begin{align*}
 I_{\eta,p}
 &=\left(h_1,\ldots,h_d,h-\sum_i z_i h_i\right)\\
 &=\left(h,h_1,\ldots,h_d\right).                       \tag{6}
\end{align*}
The second equality uses only $\sum z_i h_i\in(h_1,\ldots,h_d)$; it is not
an Euler or quasihomogeneity argument.

The scheme in (6) is the first-jet singularity scheme of the section
$X\cap W$.  This is the hypersurface specialization of the jet-incidence
description of Aluffi--Cukierman~\cite[\S1, pp.~247--249]{AluffiCukierman1993}.
We use (6) as a local calculation, not as a new general theorem about
discriminants.

Because $\gamma_X$ is finite, the section singularity is isolated.  Passing
between the analytic local ring and its completion preserves finite
colength.  Thus (6) proves the first two identities of
Theorem~\ref{thm:local}.  The third is Dimca's classical
multiplicity--Milnor formula~\cite{Dimca1986}; see also
Parusi\'nski~\cite{Parusinski1991} and the modern use in
Dimca--Ilardi~\cite[(1.3)]{DimcaIlardi2024}.  Since
\[
 (\partial h)\subset(h,\partial h),
\]
one has $\tau(h)\le\mu(h)$.  Equality is equivalent to
$h\in(\partial h)$, and Saito's theorem identifies this with analytic
quasihomogeneity~\cite{Saito1971}.  This proves (1).

To prove (2), differentiate (5).  The Hessian description of Gauss
ramification is classical; compare Piene~\cite{Piene1978}.  Here the exact
Fitting equality follows directly from the local presentation.  The first
$d$ rows are the Hessian matrix
$\operatorname{Hess}h$, while the last row is
\[
 d\!\left(h-\sum_i z_i h_i\right)=-\sum_i z_i\,dh_i.
\]
It is therefore the $(-z_1,\ldots,-z_d)$-linear combination of the first
rows.  The maximal minors presenting $\Omega_{X/X^\vee,p}$ generate exactly
$(\det\operatorname{Hess}h)$.

\begin{remark}
The fibre length is $\length(R/I_{\eta,p})=\tau(h)$.  It is not the
Hilbert--Samuel multiplicity $e(I_{\eta,p},R)$, and neither invariant should
be identified merely from the common ideal of the fibre.  In particular, the
global multiplicity--Milnor formula does not turn fibre length into Milnor
number.
\end{remark}

\section{Absorption creates infinitesimal fibre thickness}

We recall the annihilator argument in the form needed here.

\begin{lemma}[Local annihilator consequence]\label{lem:annihilator}
If $Z$ is point-span $s$-osculating-absorbing and
$a\in H^0(X,\cI_Z\otimes H^m)$, then
\[
 a_p\in\m_{X,p}^{s+1}H_p^m\qquad(p\in Z).
\]
\end{lemma}

\begin{proof}
The annihilator of $S_Z$ is $H^0(X,\cI_Z\otimes H^m)$.  The annihilator of
$\widehat\Osc^s_p(H^m)$ consists of sections whose germs vanish modulo
$\m_{X,p}^{s+1}$.  The defining containment of osculating blocks reverses
under annihilators and gives the assertion.
\end{proof}

Let $\ell_W$ be a linear equation of $W$.  For each $p\in Z$, choose a
section $g_p\in H^0(X,H^{m-1})$ nonzero at $p$, with $g_p=1$ when $m=1$.
The section $\ell_Wg_p$ of $H^m$ vanishes on all of $Z$.  Lemma
\ref{lem:annihilator} and division by the unit $(g_p)_p$ give
\[
 (\ell_W|_X)_p\in\m_{X,p}^{s+1}.                        \tag{7}
\]
In the graph coordinates of Section~\ref{sec:local}, (7) is precisely
$h_p\in\m_p^{s+1}$.  Hence
\[
 I_{\eta,p}=(h_p,\partial h_p)\subset\m_p^s.             \tag{8}
\]
The reversed ideal inclusion in (8) gives the asserted closed immersion
$\operatorname{Spec}(\cO_{X,p}/\m_p^s)\hookrightarrow\Gamma_\eta$ and
\[
 \tau(h_p)\ge\length(\cO_{X,p}/\m_p^s)
 =\binom{d+s-1}{d}.                                      \tag{9}
\]

Order-$s$ absorption contains tangent absorption.  The exact support floor
from~\cite[Theorem~1.2]{ErikssonDual2026}, ultimately based on the
point-span rank theorem~\cite{ErikssonTangent2026}, is
\[
 |Z|\ge N_{d,m}.
\]
Summing (9) proves (3).  Finally, every Hessian entry belongs to
$\m_p^{s-1}$, so its determinant belongs to $\m_p^{d(s-1)}$.  The determinant
is not the zero germ because the finite birational Gauss map is generically
separable over $\kk$.  If $s\ge2$, (9) gives local fibre length greater than
one at every support.  This completes the proof of
Theorem~\ref{thm:floor}.

\section{A defect-one isolated germ}\label{sec:defect}

The next calculation supplies the two analytic types used in the global
construction.

\begin{lemma}[Fermat and defect-one germs]\label{lem:defect}
Let $s\ge1$ and
\[
 f_0=\sum_{i=1}^d x_i^{s+1}.
\]
Then $\mu(f_0)=\tau(f_0)=s^d$.  If $d(s-1)>s+1$ and $\lambda\in\kk^*$, put
\[
 f_1=\sum_{i=1}^d x_i^{s+1}
      +\lambda\prod_{i=1}^d x_i^{s-1}.                  \tag{10}
\]
Then
\[
 \mu(f_1)=s^d,
 \qquad \tau(f_1)=s^d-1.                                \tag{11}
\]
For both germs the Hessian determinant has order exactly $d(s-1)$.
\end{lemma}

\begin{proof}
For $f_0$ the gradient ideal is $(x_1^s,\ldots,x_d^s)$, and Euler's formula
puts $f_0$ in that ideal.  This gives the first assertion.

Put $r=\prod_i x_i^{s-1}$ and $k=d(s-1)$.  The hypothesis $k>s+1$ means
that (10) has Fermat initial form.  The initial forms of its partial
derivatives are $(s+1)x_i^s$ and form a homogeneous regular sequence.
Therefore the associated graded Milnor algebra has basis
\[
 x_1^{a_1}\cdots x_d^{a_d},\qquad 0\le a_i\le s-1,
\]
and $\mu(f_1)=s^d$.  The class of $r$ is nonzero and spans the socle.  In the
Milnor algebra, Euler's identity reads
\[
 0=\sum_i x_i\partial_i f_1=(s+1)[f_0]+k\lambda[r],
\]
so
\[
 [f_1]=\left(1-\frac{k}{s+1}\right)\lambda[r]\ne0.
\]
It spans a one-dimensional ideal.  Quotienting the Milnor algebra by this
class proves $\tau(f_1)=s^d-1$.  The leading Hessian determinant in both
cases is a nonzero scalar multiple of $r$, proving the final assertion.
\end{proof}

For the family (10), the plane case first satisfies
$2(s-1)>s+1$ at $s=4$.  Its corresponding fixture is
\[
 f_1=x^5+y^5+x^3y^3,\qquad \mu(f_1)=16,\quad\tau(f_1)=15. \tag{12}
\]
It is included in the exact replay accompanying the paper.  Equation (12)
is also a counterexample to the stronger but false bound $\tau\ge s^d$.

\section{Simultaneous analytic types from prescribed jets}

We give the construction in enough detail to isolate its degree bookkeeping.
Its interpolation and extension architecture is adapted from the preceding
proper-span construction~\cite{ErikssonOsculating2026}.  Compare also the
squared-point independence criterion of
Ballico--Chiantini~\cite[Theorem~3.5]{BallicoChiantini2021}.  The general
fact that isolated hypersurface singularities can be realized in a
hyperplane section of a smooth hypersurface is classical
\cite[Proposition~11.6]{Dimca1987}; what is controlled below is the exact
support, absorption, two simultaneous analytic types, defect ladder, and an
explicit sufficient degree.

\begin{lemma}[Simultaneous jets from separators]\label{lem:jets}
Let $p_1,\ldots,p_t$ be distinct points of $\PP^d$ and let $a_i\ge0$.  If
\[
 k\ge\sum_{i=1}^t(a_i+1)-1,
\]
then restriction is surjective:
\[
 H^0(\PP^d,\cO(k))\longrightarrow
 \bigoplus_{i=1}^t
 H^0\!\left(\cO(k)\otimes\cO_{\PP^d,p_i}/\m_{p_i}^{a_i+1}\right).
\]
\end{lemma}

\begin{proof}
For $i\ne j$, choose a linear form $\ell_{ij}$ vanishing at $p_j$ and
nonzero at $p_i$.  The product
\[
 P_i=\prod_{j\ne i}\ell_{ij}^{a_j+1}
\]
kills every other target and is a unit at $p_i$.  Forms of degree $a_i$
realize all order-$a_i$ jets at $p_i$; multiplying them by $P_i$ isolates
that target.  The common degree is at most $\sum_j(a_j+1)-1$.
Multiplication by a form nonzero at all supports raises the degree.
\end{proof}

Choose $N=N_{d,m}$ distinct points $Z\subset W\cong\PP^d$ such that
\[
 H^0(W,\cO_W(m))\xrightarrow{\sim}H^0(Z,\cO_Z(m)).       \tag{13}
\]
Such a set exists because the evaluation lines of all points span the dual
of $H^0(W,\cO_W(m))$; choose a basis among them.  Set
\[
a=s^d+1,\qquad E=(a+1)N=(s^d+2)N.
\]
The binomial theorem gives $d(s-1)\le s^d-1<a$, so the perturbation term in
$f_1$ is visible in the prescribed $a$-jet.
At each point prescribe the $a$-jet of either $f_0$ or $f_1$ from Lemma
\ref{lem:defect}, choosing $f_1$ at exactly $q$ supports.  Lemma
\ref{lem:jets} supplies a section $f_*\in H^0(W,\cO_W(D))$ with those jets.
Let
\[
 V=H^0(W,\cI_{(a+1)Z}(D)),\qquad \mathcal A=f_*+V.
\]

\begin{proposition}[A section with exactly the prescribed singularities]\label{prop:section}
If $D\ge E+1$, there is $f\in\mathcal A$ such that $V(f)\subset W$ is
smooth away from $Z$ and has at the supports precisely the prescribed
contact-equivalence classes.
\end{proposition}

\begin{proof}
Greuel--Lossen--Shustin prove that an isolated germ is contact
$(\tau+1)$-determined~\cite[Corollary~2.24]{GreuelLossenShustin2007}.
Both germs have $\tau\le s^d$, so their prescribed $a=s^d+1$ jets determine
their contact classes.

Fix $u\in W\setminus Z$.  For every $p\in Z$, choose a hyperplane
$\ell_{p,u}$ through $p$ but not $u$.  The section
\[
 A_u=\prod_{p\in Z}\ell_{p,u}^{a+1}
\]
has degree $E$, belongs to $\cI_{(a+1)Z}$, and is a unit on the first
neighbourhood of $u$.  Since $D-E\ge1$, multiplication by $A_u$ shows that
$V$ surjects onto the first jets at $u$.  Thus the affine condition
$j^1_u f=0$ has codimension $d+1$ in $\mathcal A$.  Its incidence over the
$d$-dimensional base $W\setminus Z$ has dimension at most
$\dim\mathcal A-1$.  A member outside its closure is smooth away from $Z$;
finite determinacy gives the claimed germs at $Z$.
\end{proof}

Embed $W=V(y)$ in $\PP^{d+1}$ and lift $f$ independently of $y$.  For
$G\in H^0(\PP^{d+1},\cO(D-1))$, put
\[
 F_G=f+yG.
\]

\begin{proposition}[Smooth extension and exact Gauss support]\label{prop:extension}
There is $G$, nonzero at every point of $Z$, such that
$X=V(F_G)$ is smooth and integral and
\[
 (\gamma_X^{-1}([W]))_{\mathrm{red}}=Z.
\]
At each $p\in Z$, the section germ $W|_X$ is contact equivalent to the
prescribed germ of $f$.
\end{proposition}

\begin{proof}
On $W\setminus Z$, smoothness follows from Proposition~\ref{prop:section}.
At $p\in Z$, one has $dF_G=G(p)\,dy$, so $G(p)\ne0$ makes $X$ smooth.
Off $W$, multiplication by the unit $y$ and the first-jet generation of
$\cO(D-1)$ show that the singularity condition has codimension $d+2$ in the
space of $G$'s over a $(d+1)$-dimensional base.  The closure of the incidence
image is therefore a proper closed subset of the space of $G$'s.  Avoid it
together with the finitely many hyperplanes $G(p)=0$.

The hypersurface sequence and
$H^1(\PP^{d+1},\cO(-D))=0$ give $H^0(X,\cO_X)=\kk$, so the smooth
hypersurface is connected and hence integral.  Its section by $W$ is $V(f)$
and has singular locus exactly $Z$.  If a point maps under the
Gauss map to $[W]$, it lies in its tangent hyperplane $W$; within $W$, this
happens exactly where $V(f)$ is singular.  This proves the exact reduced
fibre identity.  Locally on $X$,
\[
 y=-f/G,
\]
which differs from $f$ by a unit and therefore has the prescribed contact
class.
\end{proof}

\section{Proof of the realization theorem}

Apply Propositions~\ref{prop:section} and~\ref{prop:extension}.  Since
$D>m$, restriction gives
\[
 H^0(\PP^{d+1},\cO(m))\xrightarrow{\sim}H^0(X,H^m).
\]
If a degree-$m$ ambient form $Q$ vanishes on $Z$, (13) says that
$Q|_W=0$, so $Q=yR$.  At every support, $y|_X=-f/G$ has order $s+1$.
Therefore every section of $H^m$ vanishing on $Z$ vanishes to order
$s+1$ at each support.  The reverse-annihilator argument of Lemma
\ref{lem:annihilator} proves order-$s$ absorption.

Equation (13) also gives
\[
 \dim S_Z=N.
\]
This span is proper because
\[
 h^0(X,H^m)=\binom{d+m+1}{d+1}>\binom{d+m}{d}=N.
\]

Each prescribed germ has Milnor number $s^d$.  Exactly $q$ have Tjurina
number $s^d-1$ and the others have Tjurina number $s^d$.  Theorem
\ref{thm:local} now gives
\begin{align*}
 \mult_{[W]}(X^\vee)&=Ns^d,\\
 \length(\Gamma_{[W]})&=(N-q)s^d+q(s^d-1)=Ns^d-q.
\end{align*}
Over $\kk$, both $\mu$ and $\tau$ are invariants of analytic contact
equivalence; see~\cite[Section~2.2]{GreuelLossenShustin2007}.  For the
Hessian assertion, the interpolation gives
$f-f_i\in\m^{a+1}$ with $a+1>s+1$, and on $X$ one has $y|_X=uf$ for the unit
$u=-1/G$.  Thus the degree-$(s+1)$ initial form of $y|_X$ is the Fermat form
of $f_i$, up to the nonzero scalar $u(0)$ and the invertible linear part of a
coordinate change.  Its initial Hessian determinant remains nonzero of
degree $d(s-1)$.  Lemma~\ref{lem:defect} now gives the claimed exact order
and completes the proof of Theorem~\ref{thm:realization}.

\section{Context, reproducibility, and limitations}

We use the scheme-theoretic jet-incidence formulation of
Aluffi--Cukierman~\cite{AluffiCukierman1993}; the multiplicity--Milnor
identity is due to Dimca~\cite{Dimca1986}, with broader results by
Parusi\'nski~\cite{Parusinski1991}; and the equality criterion
$\mu=\tau$ is Saito's theorem~\cite{Saito1971}.  The point-span floor and the
prescribed-jet extension architecture are authorial antecedents
\cite{ErikssonTangent2026,ErikssonOsculating2026,ErikssonDual2026}, not
independent validation.  The contribution claimed here is limited to the
exact combination of the absorption-forced fat-fibre and ramification bounds
with the explicit simultaneous defect realization at the stated sufficient
degree.

The accompanying script uses exact rational arithmetic to compute local
Macaulay quotients for the defect-one fixtures $(d,s)=(2,4),(3,3),(4,2)$,
verify the predicted Milnor and Tjurina numbers, check the Euler and initial
Hessian terms in (12), and audit finite grids of the displayed integer
formulae.  It does not mechanize analytic coordinates, finite determinacy,
incidence dimension counts, or literature priority.

The limitations are material:
\begin{itemize}
\item Everything is over $\kk$.  No positive-characteristic extension of
      the Milnor, Saito, or separability arguments is asserted.
\item The fibre is the complete scheme-theoretic fibre of the Gauss
      normalization over one reduced dual point.  Arbitrary subsets of its
      support are not covered.
\item The lower length floor (3) is not proved sharp for $s>1$.  The
      realization theorem controls a different, near-Milnor range of lengths.
\item Formula (2) describes the Fitting ramification scheme on the normal
      source.  No identification of its image with a canonical branch or
      discriminant scheme on the nonnormal dual is claimed.
\item The degree threshold (4) is deliberately sufficient and may be far
      from minimal.  No classification of equality cases or exhaustive
      priority review is claimed.
\item The computational replay certifies only its finite fixtures and
      arithmetic grids.  It is not formal proof or human peer review.
\end{itemize}

\section{Conclusion}

Osculating absorption thickens a special Gauss fibre scheme, not merely its
reduced support.  The thickness is measured locally by Tjurina algebras,
while the singularity of the dual is measured by Milnor numbers.  Their
difference is exactly the total quasihomogeneity defect of the tangent
section singularities.  The resulting dictionary
\[
 \text{fat Gauss fibre}\longleftrightarrow\sum\tau_p,
 \qquad
 \text{dual multiplicity}\longleftrightarrow\sum\mu_p
\]
both prevents a common numerical conflation and yields the global floor
proved here under the specific absorption hypothesis.  Prescribed finite
jets show that the two invariants can be separated one unit at a time across
an extremal large fibre.  The optimal Tjurina floor and the minimal ambient
degree are not determined here.

\begin{thebibliography}{99}

\bibitem{AluffiCukierman1993}
P.~Aluffi and F.~Cukierman,
\emph{Multiplicities of discriminants},
Manuscripta Math. \textbf{78} (1993), 245--258.
\url{https://doi.org/10.1007/BF02599311}

\bibitem{BallicoChiantini2021}
E.~Ballico and L.~Chiantini,
\emph{On the Terracini locus of projective varieties},
Milan J. Math. \textbf{89} (2021), 1--17.
\url{https://doi.org/10.1007/s00032-020-00324-5}

\bibitem{Dimca1986}
A.~Dimca,
\emph{Milnor numbers and multiplicities of dual varieties},
Rev. Roumaine Math. Pures Appl. \textbf{31} (1986), no.~6, 535--538.

\bibitem{Dimca1987}
A.~Dimca,
\emph{Topics on Real and Complex Singularities},
Advanced Lectures in Mathematics, Friedr. Vieweg \& Sohn, Braunschweig, 1987.
\href{https://doi.org/10.1007/978-3-663-13903-4}{doi:10.1007/978-3-663-13903-4}

\bibitem{DimcaIlardi2024}
A.~Dimca and G.~Ilardi,
\emph{On the duals of smooth projective complex hypersurfaces},
Publ. Mat. \textbf{68} (2024), 431--438.
\url{https://doi.org/10.5565/PUBLMAT6822404}

\bibitem{ErikssonTangent2026}
L.~Eriksson,
\emph{The exact rank floor for point-span tangent absorption in arbitrary
characteristic}, ARR-2026-2MHNZRRJP49Y9SWP, v3 (2026).
\url{https://arr-research.github.io/papers/ARR-2026-2MHNZRRJP49Y9SWP/versions/v3/}

\bibitem{ErikssonOsculating2026}
L.~Eriksson,
\emph{Exact floors and proper-span extremizers for higher osculating
absorption}, ARR-2026-66Q8M61AA196T8BC, v2 (2026).
\url{https://arr-research.github.io/papers/ARR-2026-66Q8M61AA196T8BC/versions/v2/}

\bibitem{ErikssonDual2026}
L.~Eriksson,
\emph{Exact multiplicity floors for dual singularities from absorbing Gauss
fibres}, ARR-2026-0WAPCGQHNC82S8VJ, v1 (2026).
\url{https://arr-research.github.io/papers/ARR-2026-0WAPCGQHNC82S8VJ/}

\bibitem{GreuelLossenShustin2007}
G.-M.~Greuel, C.~Lossen, and E.~Shustin,
\emph{Introduction to Singularities and Deformations},
Springer Monographs in Mathematics, Springer, 2007.
\url{https://doi.org/10.1007/3-540-28419-2}

\bibitem{Parusinski1991}
A.~Parusi\'nski,
\emph{Multiplicity of the dual variety},
Bull. London Math. Soc. \textbf{23} (1991), no.~5, 429--436.
\url{https://doi.org/10.1112/blms/23.5.429}

\bibitem{Piene1978}
R.~Piene,
\emph{Polar classes of singular varieties},
Ann. Sci. \`Ecole Norm. Sup. (4) \textbf{11} (1978), no.~2, 247--276.
\url{https://doi.org/10.24033/asens.1346}

\bibitem{Saito1971}
K.~Saito,
\emph{Quasihomogene isolierte Singularit\"aten von Hyperfl\"achen},
Invent. Math. \textbf{14} (1971), 123--142.
\url{https://doi.org/10.1007/BF01405360}

\bibitem{Tevelev2005}
E.~A.~Tevelev,
\emph{Projective Duality and Homogeneous Spaces},
Encyclopaedia of Mathematical Sciences 133, Springer, 2005.
\url{https://doi.org/10.1007/b138367}

\end{thebibliography}

\end{document}
